\begin{textsample}{POS Dim 3 – ap – Score 14.00 – ap/ap\_inf\_7.txt}  \label{ex:f3_pos_144}
3.3.
Eixos estruturantes das práticas pedagógicas : interações e \textbf{brincadeiras} De acordo com as Diretrizes Curriculares Nacionais de Educação \textbf{Infantil}, em seu Art. 4º, definem a \textbf{criança} como “ sujeito histórico e de direitos, que interage, \textbf{brinca}, imagina, fantasia, deseja, aprende, observa, experimenta, narra, questiona e constrói sentidos sobre a natureza e a sociedade, produzindo cultura ” ( BRASIL, 2009, p. 2 ) seres que, em suas ações e interações com os outros e com o mundo físico, constroem e se apropriam de conhecimentos.
Ainda de acordo com as DCNEI, em seu Art. 9º, os eixos estruturantes das práticas pedagógicas dessa etapa da Educação Básica são as interações e as \textbf{brincadeiras}, experiências por meio das quais as \textbf{crianças} podem construir e se apropriar de conhecimentos por meio de suas ações e interações com seus pares e com os \textbf{adultos}, o que possibilita aprendizagens, desenvolvimento e socialização.
A interação durante o \textbf{brincar} caracteriza o cotidiano da infância, trazendo consigo muitas aprendizagens e potenciais para o desenvolvimento integral das \textbf{crianças}.
Para que o \textbf{brincar} e o aprender se efetivem como práticas indissociáveis nesse nível de ensino, outro aspecto fundamental a ser considerado é a organização dos tempos e espaços de ação das \textbf{crianças}.
Cabe ao educador constituir, com elas, um ambiente rico e instigante, que caracterize a identidade do grupo, onde sejam permitidos e estimulados movimentos e organizações flexíveis, propiciando interação, questionamentos, experimentação, aprendizagem e, sobretudo, prazer, elemento propulsor do fazer \textbf{infantil}, a expressão dos \textbf{afetos}, a mediação das frustrações, a resolução de conflitos e a regulação das emoções. ( ZABALA, 1998 ).
Na Educação \textbf{Infantil} se busca assegurar as condições para que as \textbf{crianças} aprendam em situações nas quais possam desempenhar um papel ativo em ambientes que as convidem a vivenciar desafios e a se sentirem provocadas a resolvê -los, nas quais possam construir significados sobre si, os outros e os mundos social e natural.
Portanto, a sala de aula é apenas um dos espaços de construção de aprendizagem para com e para a \textbf{criança}.
De acordo com as especificidades de cada instituição, é necessário adequar diferentes atividades a espaços físicos diversos.
A construção do conhecimento é o foco das principais concepções que regem o trabalho pedagógico nas instituições de Educação \textbf{Infantil}, tornando -se necessária a criação de um ambiente acolhedor e ao mesmo tempo estimulante, em que a \textbf{criança} se sinta segura para experimentar, formular hipóteses, criar e expressar -se com liberdade. ( BRITO, 2005 ).
Assim, percebe -se a ludicidade como um importante instrumento de aprendizagem, abrindo caminho para a autonomia, a criatividade, a exploração dos significados e sentidos, além de favorecer o \textbf{equilíbrio} afetivo da \textbf{criança}.
Para tanto, acreditamos que não basta apenas à escola possuir infraestrutura adequada, mas faz -se necessário definir a concepção de infância que a escola adota, por entendermos como aspecto essencial de organização e planejamento educativo.
Os conceitos de infância e Educação \textbf{Infantil} são impregnados de história, representações, imaginação fértil por intermédio do \textbf{brincar} como ferramenta natural de desenvolvimento, valores modificam -se ao longo dos tempos, expressando o pensamento da sociedade em que foram constituídos.
Compreender a evolução histórico-cultural do conceito de infância é imprescindível aos educadores que atuam, direta ou indiretamente, com \textbf{crianças} e que buscam entender as especificidades dessa fase da vida.
Para a escola na sua concepção de infância, a \textbf{criança} deve ser pensada como um ser simultaneamente singular e social, a quem não se devem negar a individualidade e a valorização do contexto social em que está inserida.
Segundo Vygotsky ( 1996 ) este é um dos motivos pelos quais a construção realizada pelos alunos não pode ser solitária, pois o ensino precisa ser visto como um processo conjunto, compartilhado, no qual a \textbf{criança}, ajudada pelo professor e por seus \textbf{colegas}, pode mostrar -se \textbf{progressivamente} autônoma na resolução de tarefas, na utilização de conceitos, na prática de determinadas iniciativas em inúmeras questões.
A capacidade e interesse em aprender, descobrir e ampliar conhecimentos são inerentes às \textbf{crianças} \textbf{pequenas}.
Para elas, em seu cotidiano, tudo é fonte de \textbf{curiosidade} e exploração.
Agem em seu entorno, selecionando informações, analisando -as, criando relações e dando -lhes diferentes sentidos.
Dessa forma, entendem e transformam a realidade ; aprendem a respeito de si, das pessoas e do mundo, crescem e constituem suas identidades pessoais.
Essa concepção de \textbf{criança} como ser que observa, questiona, levanta hipóteses, conclui, faz julgamentos, assimila valores, constrói conhecimentos e se apropria do conhecimento sistematizado por meio da ação e nas interações com os mundos físico e social não deve resultar no confinamento dessas aprendizagens a um processo de desenvolvimento natural ou espontâneo.
Ao contrário, reitera a importância e necessidade de imprimir intencionalidade educativa às práticas pedagógicas na Educação \textbf{Infantil}, tanto na \textbf{creche} quanto na \textbf{pré-escola}. ( OSTETO, 2008 ).
A intencionalidade do processo educativo pressupõe o monitoramento das práticas pedagógicas e o \textbf{acompanhamento} da aprendizagem e do desenvolvimento das \textbf{crianças}.
O monitoramento das práticas pedagógicas se fundamenta na observação sistemática, pelo educador, dos efeitos e resultados de suas ações para as aprendizagens e o desenvolvimento das \textbf{crianças}, a fim de aperfeiçoar ou corrigir suas práticas, quando for o caso. ( BASSEDAS, 1998 ).
O \textbf{acompanhamento} da aprendizagem e do desenvolvimento se dá pela observação da trajetória de cada \textbf{criança} e de todo o grupo, levando em consideração suas conquistas, avanços, possibilidades e aprendizagens.
O \textbf{brincar} é reconhecido como uma importante linguagem que permite às \textbf{crianças} compartilhar os significados da cultura e construir sua identidade social e pessoal, é fundamental, numa instituição educativa, que ele constitua uma das formas de mediação das relações estabelecidas com as \textbf{crianças} e delas com outros sujeitos e com os objetos.
Enfim, ao compreender a importância do \textbf{brincar} para as \textbf{crianças}, o(a) professor(a) de Educação \textbf{Infantil} tem o importante papel de favorecer que ele aconteça, de forma bastante rica, no cotidiano de sua prática pedagógica.

% matched lemmas: acompanhamento, adulto, afeto, brincadeira, brincar, colega, creche, criança, curiosidade, equilíbrio, infantil, pequeno, progressivamente, pré-escola
\end{textsample}
